\chapter{Algorithms Without Computers}

\section{Introduction}

An \textbf{algorithm} is a series of steps to solve a problem.  
You do not need a computer to use algorithms—everyday tasks follow algorithms.  

Understanding algorithms helps students think logically and solve problems efficiently.

\section{Real-Life Student Problem: Making Tea}

\subsection{Case Study: Rafi’s Morning Tea}

Rafi wants to make a cup of tea. If he does not follow the steps carefully, the tea may be too weak or too strong, or the sugar may spill.  

\begin{definition}{Algorithm}
An \textbf{algorithm} is a sequence of instructions that solve a problem or complete a task.
\end{definition}

\subsection{Step-by-Step Instructions for Tea}

\begin{enumerate}
    \item Boil water in a kettle.
    \item Place tea leaves or tea bag into a cup.
    \item Pour hot water into the cup.
    \item Add sugar and milk (optional).
    \item Stir and serve.
\end{enumerate}

Following these steps ensures a perfect cup of tea every time.  

\section{Diagram: Tea-Making Algorithm}

\begin{center}
\begin{tikzpicture}[node distance=1.5cm]
\node (start) [draw, ellipse, fill=green!20] {Start};
\node (boil) [draw, rectangle, fill=yellow!20, below of=start] {Boil Water};
\node (tea) [draw, rectangle, fill=orange!20, below of=boil] {Add Tea Leaves};
\node (pour) [draw, rectangle, fill=orange!20, below of=tea] {Pour Hot Water};
\node (sugar) [draw, rectangle, fill=yellow!20, below of=pour] {Add Sugar/Milk};
\node (stir) [draw, rectangle, fill=yellow!20, below of=sugar] {Stir};
\node (serve) [draw, ellipse, fill=green!20, below of=stir] {Serve};
\draw[->] (start) -- (boil);
\draw[->] (boil) -- (tea);
\draw[->] (tea) -- (pour);
\draw[->] (pour) -- (sugar);
\draw[->] (sugar) -- (stir);
\draw[->] (stir) -- (serve);
\end{tikzpicture}
\end{center}

\section{Algorithms for Daily Life}

Students use algorithms every day, often without realizing it:

\begin{itemize}
    \item Getting ready for school
    \item Packing a bag
    \item Following a recipe
    \item Solving a math problem step-by-step
\end{itemize}

\subsection{Observation}

Algorithms help reduce mistakes, save time, and make tasks predictable.

\section{Algorithm Characteristics}

\begin{itemize}
    \item \textbf{Definiteness:} Each step is clear.
    \item \textbf{Input:} What information is required.
    \item \textbf{Output:} The result of completing the steps.
    \item \textbf{Finiteness:} The algorithm ends after a finite number of steps.
\end{itemize}

\section{Algorithm Example: Packing a School Bag}

\begin{enumerate}
    \item Check timetable.
    \item Collect books for today’s subjects.
    \item Collect homework notebooks.
    \item Pack stationery.
    \item Pack water bottle and lunch.
    \item Close the bag and carry to school.
\end{enumerate}

\begin{definition}{Daily Life Algorithm}
Any routine task can be written as an algorithm using steps that anyone can follow.
\end{definition}

\section{Step-by-Step Diagram: Packing Bag}

\begin{center}
\begin{tikzpicture}[node distance=1.5cm]
\node (start) [draw, ellipse, fill=green!20] {Start};
\node (timetable) [draw, rectangle, fill=yellow!20, below of=start] {Check Timetable};
\node (books) [draw, rectangle, fill=orange!20, below of=timetable] {Collect Books};
\node (homework) [draw, rectangle, fill=orange!20, below of=books] {Collect Homework};
\node (stationery) [draw, rectangle, fill=yellow!20, below of=homework] {Pack Stationery};
\node (lunch) [draw, rectangle, fill=yellow!20, below of=stationery] {Pack Lunch \& Water};
\node (done) [draw, ellipse, fill=green!20, below of=lunch] {Done};
\draw[->] (start) -- (timetable);
\draw[->] (timetable) -- (books);
\draw[->] (books) -- (homework);
\draw[->] (homework) -- (stationery);
\draw[->] (stationery) -- (lunch);
\draw[->] (lunch) -- (done);
\end{tikzpicture}
\end{center}

\section{Key Concepts}

\begin{itemize}
    \item An algorithm is a set of instructions to solve a problem.
    \item Algorithms are part of everyday life.
    \item Algorithms must be definite, finite, and produce a result.
    \item Writing steps clearly reduces mistakes.
\end{itemize}

\section{End-of-Chapter Exercises}

\begin{exercise}{Very Short Questions}
\begin{enumerate}
    \item What is an algorithm?
    \item Give two examples of algorithms in daily life.
    \item What is meant by definiteness in an algorithm?
\end{enumerate}
\end{exercise}

\begin{exercise}{Short Answer Questions}
\begin{enumerate}
    \item Write the algorithm for brushing teeth.
    \item Explain why packing a school bag is an algorithm.
    \item How does following an algorithm reduce mistakes?
\end{enumerate}
\end{exercise}

\begin{exercise}{Activity-Based Questions}
\begin{enumerate}
    \item Create an algorithm for making your favorite sandwich.
    \item Observe a classmate performing a routine task and write an algorithm for it.
    \item Identify input, process, and output for a daily task at home.
\end{enumerate}
\end{exercise}

\section{Teacher Notes}

\marginnote{
\textbf{Learning Objectives:} 
\begin{itemize}
    \item Students understand what an algorithm is and how it applies to real life.
    \item Students can write step-by-step instructions for daily tasks.
    \item Students recognize input, process, and output in everyday algorithms.
\end{itemize}

\textbf{Teaching Tips:} 
\begin{itemize}
    \item Demonstrate a real-life task in class and write steps together.
    \item Encourage students to identify steps for complex daily routines.
\end{itemize}

\textbf{Assessment:} 
\begin{itemize}
    \item Ask students to write their own daily-life algorithms.
    \item Check clarity and completeness of steps.
\end{itemize}
}